% Figure: a block on a rough inclined plane, real scene (left) and
% free-body diagram (right). The heart of the chapter: turning a scene
% into vectors.
\begin{figure}[htbp]
\centering
\begin{tikzpicture}[
  weight/.style={draw=engblue, -{Stealth}, very thick},
  normal/.style={draw=engred, -{Stealth}, very thick},
  fric/.style={draw=engamber, -{Stealth}, very thick},
  axis/.style={draw=enggray, -{Stealth}, thin, dashed},
  scale=1.0,
]
% === Left: the scene ===
\begin{scope}[xshift=0cm]
  % incline triangle, angle ~ 25 deg
  \coordinate (A) at (0,0);
  \coordinate (B) at (5.2,0);
  \coordinate (C) at (5.2,2.42);
  \draw[draw=engblack, very thick] (A) -- (B) -- (C) -- cycle;
  % hatching on ground
  \foreach \x in {0.3,0.8,...,4.8} {
    \draw[draw=enggray, thin] (\x,0) -- (\x-0.25,-0.25);
  }
  % angle
  \draw[draw=enggray, thin] (1.1,0) arc (0:25:1.1);
  \node[font=\scriptsize] at (1.4,0.22) {$\theta$};
  % block on the incline
  \begin{scope}[shift={(3.3,1.54)}, rotate=25]
    \draw[draw=engblack, fill=engblue!12, thick] (-0.55,0) rectangle (0.55,0.85);
    \node[font=\scriptsize] at (0,0.42) {$m$};
  \end{scope}
  \node[font=\small, anchor=north] at (2.6,-0.7) {(a) the scene};
\end{scope}
% === Right: the FBD of the block ===
\begin{scope}[xshift=8.3cm, yshift=0.6cm]
  % the block, tilted 25 deg, centered at origin
  \begin{scope}[rotate=25]
    \draw[draw=engblack, fill=engblue!8, thick] (-0.7,-0.45) rectangle (0.7,0.45);
  \end{scope}
  \coordinate (G) at (0,0);
  % weight straight down
  \draw[weight] (G) -- ++(0,-2.3) node[anchor=north, font=\small, engblue] {$mg$};
  % normal perpendicular to incline (incline at 25 deg, normal at 25+90=115 deg)
  \draw[normal] (G) -- ++(115:2.2) node[anchor=south, font=\small, engred] {$N$};
  % friction up the incline (incline direction 25 deg, up-slope = 25+180=205 deg)
  \draw[fric] (G) -- ++(205:2.0) node[anchor=east, font=\small, engamber] {$F_f$};
  % local axes along/normal to incline
  \draw[axis] (G) -- ++(25:1.4) node[anchor=west, font=\scriptsize, enggray] {$x$ (up-slope $+$)};
  \node[font=\small, anchor=north] at (0,-2.9) {(b) free-body diagram};
\end{scope}
\end{tikzpicture}
\caption[A block on a rough incline and its free-body diagram]{The
move at the centre of the chapter. (a) A block of mass $m$ rests on a
rough plane inclined at $\theta$. (b) Isolate the block and replace
every contact with the force it exerts: the weight $mg$ acts straight
down at the centre of mass, the normal force $N$ acts perpendicular to
the incline surface, and the friction force $F_f$ acts along the
surface, here drawn up-slope to oppose impending downhill slip. The
basis is chosen along and perpendicular to the incline so that only the
weight needs resolving: $\sum F_x: F_f - mg\sin\theta = 0$ and $\sum
F_y: N - mg\cos\theta = 0$. No part of the incline appears in the
diagram; only the forces it transmits.}
\label{fig:vol03ch01:fbd-block-incline}
\end{figure}
